\begin{prob}
    Consider the following code which constructs a numpy array of $n$ random numbers:\footnote{Note that in practice
        you wouldn't do this with a loop; you'd write \python{np.random.uniform(n)} to generate
    the array in one line of code.}

    \begin{minted}[autogobble]{python}
        import numpy as np
        results = np.array([])
        for i in range(n):
            results = np.append(results, np.random.uniform())
    \end{minted}

    Remember that we have to write \python{results = np.append(results, np.random.uniform())}
    instead of just \python{np.append(results, np.random.uniform())} because it turns out that
    \python{np.append} returns a \emph{copy} of \python{results} with the new
    entry appended to the end.

    Note that this code is very similar to how we taught you to run simulations in
    DSC 10: we first created an empty numpy array, and then ran our simulation in a
    loop, appending the result of each simulation with \python{np.append}.
    When we ran simulations, we often used $n = 100{,}000$ or larger (and they took a while
    to finish).

    \begin{subprobset}

        \begin{subprob}
            Guess the time complexity of the above code as a function of $n$. Don't worry
            about getting the right answer (we won't grade for correctness). You don't
            need to explain your answer.

            \begin{soln}
                I don't know... $\Theta(n)$?
            \end{soln}

        \end{subprob}

        \begin{subprob}
            Time how long the above code takes when $n$ is: $10{,}000$, $20{,}000$, $40{,}000$, $80{,}000$, $120{,}000$, and $160{,}000$.
            Then make a plot of the times, where the $x$-axis is $n$ (the input size) and the $y$-axis is the time taken in seconds.

            Remember to provide not only your plot, but to show your work by providing the code that generated it.

            Hint: You can do the timing by hand with the \%\%\python{time} magic function in a Jupyter notebook, or you can use the \python{time()} function in the \python{time}
            module. For example, to time the function \python{foo}:

            \begin{minted}[autogobble]{python}
                import time
                start = time.time()
                foo()
                stop = time.time()
                time_taken = stop - start
            \end{minted}

            \begin{soln}
                You should get a plot that looks like this:

                \includegraphics[width=\textwidth]{\thisdir/include-solution/times.pdf}

                This plot was generated by the following code:

                \inputminted{python}{\thisdir/include-solution/code.py}

            \end{soln}

        \end{subprob}

        \begin{subprob}
            Looking at your plot, what do you now think the time complexity is? Why does the code have this time complexity?

            Hint: what is the time complexity of \python{np.append}, and why?

            \begin{soln}
                It takes $\Theta(n^2)$ time (quadratic).

                This is because \python{np.append} takes linear time since it makes a copy of the array. On the first iteration
                of the loop, the array is empty and copying it takes no time. One the second iteration, one element has to be
                copied. On the third iteration, two elements have to be copied... and so on. On the $n$th iteration, $n-1$ elements
                have to be copied -- which might take some time!

                This is therefore an example of a dependant nested loop where the inner loop body takes $\Theta(i)$ time. This gives an
                arithmetic sum for the time complexity
            \end{soln}

        \end{subprob}

        \begin{subprob}
            It turns out that creating an empty numpy array and appending to it at the end of each iteration is a \emph{terrible}
            way to do things, and you should \emph{never} write code like this if you can avoid it.\footnote{We taught you the \python{np.append} way because it was conceptually simpler -- we didn't need to introduce you to Python lists. This is one instance in which your professors lie to you in early courses, then correct their lies later on.} Instead, you should create an
            empty Python list, append to it, then make an array from that list, like so:

            Remember to provide not only your plot, but to show your work by providing the code that generated it.

            \begin{minted}[autogobble]{python}
                lst = []
                for i in range(n):
                    lst.append(np.random.uniform())
                arr = np.array(lst)
            \end{minted}

            To check this, repeat part (b), but with this new code. Show your
            plot. It is OK if your plot is a little odd, but it shouldn't be
            quadratic! (Check with a tutor if you're concerned).

            \begin{soln}
                You should get something like this:

                \includegraphics[width=\textwidth]{\thisdir/include-solution/times-bar.pdf}

                It turns out that this approach has linear time complexity (to be specific, it has \emph{amortized} linear time complexity).
            \end{soln}

        \end{subprob}

    \end{subprobset}

\end{prob}
